\documentclass[10pt,a4paper]{amsart}
\usepackage[margin=19mm]{geometry}
\usepackage{amsmath,amssymb,mathtools}
\usepackage[scr=rsfs,cal=euler]{mathalfa}
\usepackage{xfrac}
\usepackage[unicode,hidelinks,bookmarks=false]{hyperref}
\newcommand{\ie}{i.e.\ }
\newcommand{\cf}{cf.\ }
\newcommand{\resp}{respectively\ }
\newcommand{\Subsection}{\subsection}
\providecommand{\nobreakdash}{\nobreak\hbox{-}}
\setlength{\emergencystretch}{2em}
\multlinegap0pt
\usepackage{siunitx}
\usepackage{siunitx}
\usepackage{siunitx}
\usepackage{siunitx}
\usepackage{amssymb}
\usepackage{tikz}
\usepackage{dsfont}
\usepackage{float}
\usepackage{siunitx}
\newtheorem{proposition}{Proposition}
\title{A Metric-Deformed $q$-Gauge Dirac Equation}
\author{Julio Cesar Jaramillo Quiceno$^a$}
\date{Source: arXiv:2605.21508v2 (CC BY 4.0)}
\begin{document}
\maketitle

\noindent\emph{Source and licence.} A Metric-Deformed $q$-Gauge Dirac Equation. Version arXiv:2605.21508v2 (CC BY 4.0), distributed by arXiv under the Creative Commons Attribution 4.0 International licence, http://creativecommons.org/licenses/by/4.0/, as recorded in the arXiv licence metadata for this exact version and shown on the abstract page https://arxiv.org/abs/2605.21508v2. Full licence notice: \texttt{licenses/arxiv-2605.21508-CC-BY-4.0.txt}.\par\smallskip

\noindent\textit{Editorial scope.} Counted unit: the article's proposition prop:YM\_gauge\_invariance (source0.tex lines 654--672). The article's complete original proof of it is delivered with it (source0.tex lines 674--724, one \texttt{proof} environment of 51 lines). No mathematics is written, repaired or re-derived: the statement and the proof are the article's own lines, in the article's own notation, and no retained line is substituted or re-flowed.  No further source-local context is needed: every cross-reference of every retained line resolves inside the two delivered spans.  Nothing was omitted from any retained span, so the package carries no omission record.  No retained line contains a citation, so the wrapper carries no bibliography and no bibliography entry.  No retained line contains a figure environment, an image-inclusion command or drawing code, so no image is delivered and the package declares no asset dependency.  Editorial formatting only, in addition: an AMS \texttt{amsart} preamble that restates the article's own declarations and macros used by the retained lines, namely the environments \texttt{proposition} on separate counters, together with the standard packages those lines need.  The retained environments print in this document as Proposition 1 in order of appearance, whereas the article prints the same items with the numbering of its own full document.  The complete native source of the article as distributed in the arXiv e-print for this exact version is bundled unchanged as \texttt{sources/201065/source0.tex}.
\begin{proposition}[Gauge invariance of the deformed Yang--Mills action]
\label{prop:YM_gauge_invariance}

The action \(S_{\mathrm{YM}}^{(q)}\) is invariant under the gauge transformations
\begin{equation}
A_a
\longrightarrow
U A_a U^{-1}
+
U(\partial_a U^{-1}),
\qquad
F_{ab}^{(q)}
\longrightarrow
U F_{ab}^{(q)} U^{-1},
\label{eq:gauge_transform_YM}
\end{equation}
where \(U(x)\in G\).

\end{proposition}

\begin{proof}

Under the gauge transformation,
\[
F_{ab}^{(q)}
\longrightarrow
U F_{ab}^{(q)} U^{-1}.
\]

Hence,
\[
F_{ab}^{(q)}F_{(q)}^{ab}
\longrightarrow
U
\left(
F_{ab}^{(q)}F_{(q)}^{ab}
\right)
U^{-1}.
\]

Using cyclicity of the trace,
\[
\operatorname{tr}(UXU^{-1})
=
\operatorname{tr}(X),
\]
we obtain
\[
\operatorname{tr}
\left(
F_{ab}^{(q)}F_{(q)}^{ab}
\right)
\longrightarrow
\operatorname{tr}
\left(
F_{ab}^{(q)}F_{(q)}^{ab}
\right).
\]

Since the effective metric is treated as a fixed background structure, the measure
\[
d^{d_{\mathrm{eff}}}x
\sqrt{
\left|
\det g_{\mathrm{eff}}
\right|
}
\]
is invariant under internal gauge transformations. Therefore the action is gauge-invariant.

\end{proof}

\end{document}
